\subsection{A central double cover and an actual mixed-prime extension}

\begin{statement}{D.48}{Proved}{The all-Fermat replacement extends over a new mixed three-primary kernel}
Keep the marked groups and the action of D.47. Every subgroup
and commutator below is closed. Put
\[
 \begin{aligned}
 L&=[\widehat N,\widehat B],&
 \mathcal E&=\widehat B/L,\\
 M^+&=L\cap\widehat F'=L\cap M,&
 \mathcal H^+_\infty&=\widehat B/M^+,\\
 L_m^+&=L\cap F_m,&
 \mathcal H_m^+&=\widehat B/L_m^+.
 \end{aligned}                                                     \tag{D.48.1}
\]
Then \(\mathcal E\to Q\) is an actual central extension with
kernel \(C_2\), \(|\mathcal E|=1152\), and the discrete subgroup
whose completion is \(L\) is free of rank \(193\).
The image of \(M\) in this central kernel is all of \(C_2\).
The image of \(L\) in the original loop abelianization is
exactly the lattice \(\Lambda\) of D.47.5. In particular
\[
 \begin{gathered}
 [M:M^+]=2,\qquad
 \mathcal H^+_\infty
    =(\widehat B/\widehat F')\times_{Q/J}\mathcal E,
       \qquad J=q(\widehat F')\simeq Q_8,\\
 \ker(\mathcal H^+_\infty\to\mathcal H_\infty)=C_2,
 \qquad
 |\mathcal H_m^+|=1152m^2/c_m,\qquad
 \operatorname{rank}(L_m^+)=1+192m^2/c_m,
 \end{gathered}                                                   \tag{D.48.2}
\]
where \(c_m=\gcd(2,m)^2\gcd(3,m)\). Every displayed map is
induced by the same original \(\widehat B\).

Define the new relative sources by
\[
 \begin{aligned}
 R_p^+&=\ker(M^+\to(M^+)^{(p)}),&
 \mathcal P_p^+&=\widehat B/R_p^+,\\
 R^+&=\bigcap_pR_p^+,&
 \mathcal V^+&=\widehat B/R^+.
 \end{aligned}                                                     \tag{D.48.3}
\]
They have exact marked descriptions
\[
 \begin{gathered}
 \mathcal P_p^+
   =\varprojlim_m\widehat B/
                \ker(L_m^+\to(L_m^+)^{(p)}),\qquad
 \mathcal V^+=\prod_{\mathcal H^+_\infty}\mathcal P_p^+,\\
 \ker(\mathcal V^+\to\mathcal H^+_\infty)
      =\prod_p(M^+)^{(p)},\qquad
 R_2^+=R_{\infty,2},\quad
 \mathcal P_2^+=\mathcal P_{\infty,2}.
 \end{gathered}                                                   \tag{D.48.4}
\]
The base \(\mathcal H^+_\infty\) is the exact maximal common
marked quotient of different prime factors and of one factor
with all the remaining ones. For odd \(p\), the actual inclusion
\(M^+\subset M\) gives an epimorphism
\[
 (M^+)^{(p)}\twoheadrightarrow M^{(p)},\qquad
 E_p=\ker((M^+)^{(p)}\to M^{(p)}).
\]
The comparisons with D.47, including both source kernels, are
\[
 \begin{gathered}
 R^+\subset R_\infty,\qquad
 1\longrightarrow\prod_{p\ne2}E_p
       \longrightarrow\mathcal V^+
       \longrightarrow\mathcal V\longrightarrow1,\\
 \ker(\mathcal V\to\mathcal H^+_\infty)
      =(M^+)^{(2)}\times\prod_{p\ne2}M^{(p)},\qquad
 \ker(\mathcal V^+\to\mathcal P_p^+)
      =\prod_{q\ne p}(M^+)^{(q)}.
 \end{gathered}                                                   \tag{D.48.5}
\]
In the second line \((M^+)^{(2)}\) is identified with the actual
index-two subgroup \(\ker(M^{(2)}\to C_2)\).

The actual arithmetic element \(h\) of D.46 fixes
\(\mathcal E\), and hence \(\mathcal H^+_\infty\), pointwise.
Consequently
\[
 \alpha^+=(\alpha_{h,\infty,2},(1_{\mathcal P_p^+})_{p\ne2})
       \in\operatorname{Aut}(\mathcal V^+)                       \tag{D.48.6}
\]
is a continuous peripheral automorphism with multiplier one.
It projects to the entire D.47 action and to its exact parameter.
One compatible pair of loop conjugators is constructed below;
the required correction to the two-primary \(r\)-conjugator
centralizes \(r\) and has trivial \(Q\)-image. Thus all the
retained finite CM data, every Fermat homology action, and the
nonarithmetical full metabelian outer class of D.47 are preserved.

This extension adds a concrete finite source which does not
factor through \(\mathcal V\). Namely let
\[
 \mathcal J_3=\widehat B/(L'L^3),\qquad
 V_3=L/(L'L^3)\simeq(\mathbb F_3)^{193},\qquad
 G_3=\widehat B/(\widehat N'\widehat N^3).
                                                               \tag{D.48.7}
\]
Here \(L^3\) is the closed subgroup generated by cubes.
The central involution \(z\in\mathcal E\) acts on \(V_3\), with
\[
 \dim V_3^+=97,\qquad\dim V_3^-=96.
                                                               \tag{D.48.8}
\]
There are exact source maps and kernel identifications
\[
 \begin{gathered}
 \mathcal P_3^+\twoheadrightarrow\mathcal J_3,\qquad
 \ker(\mathcal J_3\to
                 \mathcal E\times_QG_3)=V_3^-,\\
 \operatorname{im}(R_\infty\to\mathcal J_3)=V_3^-
                      \simeq C_3^{96},\qquad
 \operatorname{im}(R^+\to\mathcal J_3)=1.
 \end{gathered}                                                   \tag{D.48.9}
\]
Thus the maximal common marked quotient of
\(\mathcal V\) and \(\mathcal J_3\) is exactly
\(\mathcal E\times_QG_3\), of order \(1152\cdot3^{97}\), whereas
\(|\mathcal J_3|=1152\cdot3^{193}\).
The new action on \(\mathcal J_3\) is identity, and it agrees
with the D.47 action on this complete common quotient.
The actual image of \(M\) in \(\mathcal J_3\) has structure
\[
 \bigl((\mathbb F_3)^{95}\oplus(\mathbb F_3)^{96}\bigr)
                 \rtimes C_2,
                                                               \tag{D.48.10}
\]
where the involution is identity on the first summand and
minus identity on the second. This gives the precise mixed
prime component killed in D.47 and retained here.

The remaining source kernel \(R^+\) still has exact image
\(A_7\) under D.46's displayed \(S_7\) source map. Full
profinite GT0 realization of the prescribed CM pair and of
the metabelian action remains undecided. The new action is
proved on the stated larger source; its finite projections
are not substituted for a full-source lift.

The manuscript passages are Part D(e), section 51.16,
folio 734 verso, scan 174, and the two-torus and Fermat
follow-ups in section 51.15, folios 725--727, scans 156--160.
\end{statement}
\begin{proof}
\emph{The actual central extension and the lattice map.}
The image of \([\widehat N,\widehat B]\) in
\(\widehat N^{\rm ab}\) is the closed augmentation image
generated by \((1-s)\widehat N^{\rm ab}\) and
\((1-r)\widehat N^{\rm ab}\). The columns
\((1-s)b+s(1-r)a\) of D.26 generate the same subgroup:
\(s v\equiv v\) modulo the first summand. D.26's exact
integral Smith form \((1^{96},2)\) therefore identifies
\(\widehat N/L=C_2\). This proves the central extension
and its order. Its finite-index torsion-free kernel has
rank \(1+1152/6=193\), by the actual Bass--Serre graph.

The program \texttt{d\_cm\_central\_mixed.py} reconstructs
this same quotient. It first forms the oriented graph of
\(N\backslash T\), its 97 integral fundamental cycles,
and the unique nonzero modulo-two coinvariant functional.
Let \(P(q)\) be the path attached to the original word
for \(q\in Q\), and extend that functional to the edge
module by projection to non-tree edges. The actual central
cocycle is
\[
 \delta(q,q')=\phi\bigl(P(q)+qP(q')-P(qq')\bigr)\in\mathbb F_2.
                                                               \tag{D.48.11}
\]
The marked elements \((q,a)\) multiply by this cocycle.
The certificate retains their full right \(s,r\) permutations,
checks the two torsion relations at all 1152 states, and
generates every state. In particular \(t^{12}\) maps to
the central generator. It also retains an explicit
88-letter word in \(F'\) mapping to that generator;
its \(Q\)-image and both integral loop exponents vanish.
Thus the word belongs to \(M\) and proves directly
that \(M\to C_2\) is onto.

There is a second integral check of the same assertion.
On the lattice basis \((2,-2),(0,6)\), the actual
\(r\)-action is
\[
                   \begin{pmatrix}1&-3\\1&-2\end{pmatrix}.
                                                               \tag{D.48.12}
\]
Its reduction modulo two has no invariant nonzero covector.
A \(\widehat B\)-invariant character
\(\widehat N\to C_2\) vanishing on \(M\) would factor
through such a covector of \(\Lambda\), which is impossible.
Since \(M\to C_2\) is onto, multiplying a preimage of any
\(\lambda\in\Lambda\) by an appropriate element of \(M\)
makes it lie in \(L\). Hence \(L\to\Lambda\) is onto.

The code computes \(|q_{\mathcal E}(F)|=192\) and
\(|q_{\mathcal E}(F')|=16\) using actual words. The kernel
of this last group over \(J\) is precisely the central
\(C_2\), so \(\mathcal E/q_{\mathcal E}(F')=Q/J\).
The two original quotient kernels now give the fibre product
in (D.48.2) by closed Goursat, with intersection \(M^+\).
Because \(M\subset L_m\) already surjects onto this \(C_2\),
every \(L_m^+\) has index two in \(L_m\). D.47's index
and graph formulas give all finite orders and ranks.

\emph{Both relative-prime kernel comparisons.}
The subgroup \(M^+\) has index two in \(M\). The induced
pro-two topology is its full pro-two topology: a normal
finite two-power quotient of \(M^+\) is replaced by the
intersection of its two \(M\)-conjugate kernels; the
resulting quotient of \(M\) is still a finite two-group.
Restriction gives the reverse inclusion of residual kernels.
Thus \(R_2^+=R_{\infty,2}\), and
\((M^+)^{(2)}=\ker(M^{(2)}\to C_2)\), with the stated maps.

For odd \(p\), the image of \(M^+\) in \(M^{(p)}\)
is closed and normal, and its quotient is both a pro-\(p\)
group and a quotient of \(C_2\). That quotient is trivial.
Consequently \((M^+)^{(p)}\to M^{(p)}\) is onto and
\(R_p^+\subset R_{\infty,p}\). This proves
\(R^+\subset R_\infty\) and the actual source map
\(\mathcal V^+\to\mathcal V\).

For the inverse limit in (D.48.4), restriction immediately
gives
\(R_p^+\subset\ker(L_m^+\to(L_m^+)^{(p)})\).
For the other inclusion, take the kernel \(C\) of a finite
\(p\)-quotient of \(M^+\). Uniform continuity of conjugation
on the compact group \(M^+\) makes its \(\widehat B\)-orbit
finite. Intersecting that orbit replaces \(C\) by a
\(\widehat B\)-normal open subgroup with finite \(p\)-group
quotient. Choose an open normal subgroup \(O\) of
\(\widehat B\) with \(O\cap M^+=C\), by separating the
finite subgroup \(M^+/C\) in \(\widehat B/C\).
The subgroups \(L_m^+/M^+\) form a decreasing divisibility
system with intersection one in \(\mathcal H^+_\infty\).
Compactness supplies \(m\) with \(L_m^+\subset OM^+\).
Its image in \(\widehat B/O\) is a \(p\)-group, so its
pro-\(p\) residual lies in \(O\). The intersection of
these residual kernels already lies in \(M^+\).
Intersection over \(m\)
and then over all such \(C\) gives the reverse inclusion.
This proves the equality of source kernels before passage
to the inverse limit.

Distinct prime quotients of \(M^+\) have no common
nontrivial finite quotient. Finite Goursat and compactness
therefore give its full product of maximal pro-prime images.
They prove the fibre product and the maximal-common-quotient
statements in (D.48.4). Under the D.47 map
\(M\to\prod_p M^{(p)}\), the image of \(M^+\) is
\((M^+)^{(2)}\times\prod_{p\ne2}M^{(p)}\): its sole
additional defining condition is the displayed two-primary
central character. Applying the same source map to the
new product proves exactly (D.48.5), including the product
of the odd-prime kernels \(E_p\).

\emph{Pointwise base invariance and compatible conjugators.}
The actual \(h\) fixes \(Q\), preserves \(\widehat N\)
and hence preserves \(L\). In \(\mathcal E\) its images
of \(s,r\) differ from \(s,r\) by central elements
\(z_s,z_r\). The relations imply \(z_s^2=z_r^3=1\).
The exact fixed-cusp equality \(\alpha_h(sr)=sr\) gives
\(z_sz_r=1\). An element of order dividing both two and
three is identity, so both deviations vanish. Thus \(h\)
fixes \(\mathcal E\) pointwise. D.47 already proves it
fixes \(\widehat B/\widehat F'\); the exact fibre product
in (D.48.2) proves pointwise invariance of
\(\mathcal H^+_\infty\). Formula (D.48.6) is consequently
an automorphism of the actual product, preserving the
cusp and the two torsion conjugacy classes.

Here is the precise conjugator correction, including its
projection to the already retained action. Write
\(A^\dagger,\beta^\dagger\) for a pair constructed in
D.47, so their base coordinates there are respectively
\(1,\beta_\infty\). The two-primary source is unchanged.
Let \(a\in C_2\) be the image of its
\(A^\dagger_2\) in the new base. The map
\(M^{(2)}\to C_2\) has pro-two kernel, and \(r\)
has order three. Coprime fixed-coset lifting supplies
\(d\in C_{M^{(2)}}(r)\) above \(a^{-1}\).
Replace \(A^\dagger_2\) by \(A^\dagger_2d\), and
take the identity \(r\)-conjugator at every odd prime.
These have the same new base coordinate, namely one.
The change centralizes \(r\), and \(q(d)=1\), so it
changes neither the action nor its finite CM coordinate.

Retain \(\beta^\dagger_2\) and denote its image in
\(\mathcal H^+_\infty\) by \(\beta_+\).
It centralizes \(s\), because the actual \(h\) is
identity on this base. At each odd \(p\), coprime
order-two fixed-coset lifting through the pro-\(p\)
kernel supplies a centralizing lift of \(\beta_+\)
in \(\mathcal P_p^+\). Its image in
\(\mathcal P_{\infty,p}\) and the already chosen
\(\beta^\dagger_p\) differ by an element of
\(C_{M^{(p)}}(s)\). The surjection
\((M^+)^{(p)}\to M^{(p)}\) has pro-\(p\) kernel,
so the same coprime fixed-element argument lifts this
difference inside \(C_{(M^+)^{(p)}}(s)\). Multiplication
by that lift makes the projection exactly
\(\beta^\dagger_p\). These corrections are made once
in the complete groups; finite complements and compactness
give the stated lifts, as in D.46--D.47.

The resulting tuples belong to the actual fibre product.
The whole \(\beta\)-conjugator projects exactly to
D.47's \(\beta^\dagger\), whereas the \(A\)-conjugator
projects to \(A^\dagger\) times the explicit
\(r\)-centralizing tuple \((d,1,1,\ldots)\).
Thus the automorphism and the parameter
\(f=s\beta s\beta^{-1}\) project exactly to D.47.
The new parameter is again \((f_{h,2},1,1,\ldots)\),
has trivial \(\mathcal H^+_\infty\)-image, and has
an actual source preimage in \(M^+\subset\widehat F'\).
The two hexagons and generation follow from the actual
automorphism, with multiplier one and no further cusp
gauge. This proves the compatibility assertions in the
statement, rather than merely an unrelated lift of a
finite CM matrix pair.

\emph{The exact mixed quotient and the newly retained kernel.}
The quotient \(\mathcal J_3\) is the actual Bass--Serre
homology extension over \(\mathcal E\), with kernel
\(V_3\). Conjugation by a lift of \(z\) is independent
of that lift on this abelian kernel and has order two.
The program builds the full 1152-edge graph, its 193
fundamental cycles and the permutation of edges that
toggles the central coordinate. It verifies the action
on every full edge chain, its square, and the two
eigenspace ranks in (D.48.8). It also computes the
actual graph pushforward \(V_3\to\Lambda/3\Lambda\)
and verifies its rank two and its invariance under \(z\).

The map \(L\subset\widehat N\) induces
\[
              (V_3)_{C_2}\xrightarrow{\ \sim\ }
                         \widehat N^{\rm ab}/3.
                                                               \tag{D.48.13}
\]
For an explicit proof, the two-group quotient in
\(1\to L\to\widehat N\to C_2\to1\) has vanishing
positive homology with \(\mathbb F_3\) coefficients:
averaging over it contracts the positive-dimensional
resolution. The homology five-term sequence therefore
gives the displayed isomorphism with its actual inclusion
map. Equivalently the full relative pro-three quotient
of \(\widehat N\) kills the order-two deck action.
Because two is invertible in \(\mathbb F_3\), the
coinvariant kernel is exactly \(V_3^-\).

The two source maps from \(\mathcal J_3\) to
\(\mathcal E\) and \(G_3\) have exact common quotient
\(Q\): their kernels over \(Q\) have orders two and
\(3^{97}\), respectively. Their joint image is the
full fibre product by finite Goursat. Equation (D.48.13)
therefore proves its kernel is \(V_3^-\), and proves
the two displayed finite orders. This quotient already
factors through \(\mathcal V\). Indeed \(\mathcal E\)
is a quotient of its two-primary factor because
\(M\to C_2\) is two-primary, and \(G_3\) is a quotient
of its three-primary factor. Their actions are both
identity: \(h\) fixes \(\mathcal E\), and the three-primary
factor has identity action. Hence
\[
              \operatorname{im}(R_\infty\to\mathcal J_3)
                              \subset V_3^- .                 \tag{D.48.14}
\]

For the reverse inclusion, the image of \(M\) in
\(\mathcal J_3\) contains an element over \(z\), and
its commutators with the image of \(L\) give all of
\((z-1)V_3=V_3^-\). This finite image has a normal
three-subgroup of index two. A complement exists by
the elementary coprime complement theorem, and its
involution acts by minus identity on \(V_3^-\).
Every nilpotent quotient must kill this subspace:
its two- and three-primary subgroups commute, while
\(z-1=-2\) is invertible on the subspace. Conversely,
after killing it the image is abelian. Thus its exact
nilpotent residual is \(V_3^-\).

The image of the full pronilpotent residual of a compact
group in a finite quotient is that finite quotient's
nilpotent residual. To verify both inclusions, every
nilpotent quotient of the finite image factors through
the original maximal pronilpotent quotient, while
every finite quotient of that pronilpotent group is
nilpotent. Apply this to \(M\to\operatorname{im}(M\to
\mathcal J_3)\). It proves equality in (D.48.14).
For the finer structure (D.48.10), note that
\(L\to\Lambda\) is onto and the image of
\(L'L^3\) there is exactly \(3\Lambda\).
The exact maximal common quotient of
\(\operatorname{im}(\widehat N\to\mathcal J_3)\)
and \(\Lambda\) is consequently
\(\Lambda/3\Lambda\). Its kernel is the actual image
of \(M\), by closed Goursat. The rank-two invariant
map removes two dimensions from \(V_3^+\), leaving
the dimensions 95 and 96 stated there.

Finally \(M^+\subset L\), so its image in
\(\mathcal J_3\) is an abelian three-group. The actual
source map therefore kills \(R_3^+\), giving
\(\mathcal P_3^+\to\mathcal J_3\) and the second
kernel image in (D.48.9). Its action is identity by
(D.48.6). We have thus extended the retained action
across the entire indicated \(C_3^{96}\) kernel,
including the complete common quotient and the
normalized conjugator correction.

The base \(\mathcal H^+_\infty\) is prosolvable and
the new kernel over it is pronilpotent. Thus
\(\mathcal V^+\) is prosolvable. Under the literal
\(S_7\) source of D.46 its full source kernel contains
\(A_7\), and it is contained in \(A_7\) because
\(R^+\subset\widehat N\subset B_{\rm even}\).
Its image is exactly \(A_7\). This proves the final
residual assertion without a full-source lifting claim.
\end{proof}
